\documentclass{article}

\usepackage{float}
\restylefloat{table}

\usepackage{booktabs}

\title{Team Productivity: POC\\\progname}

\author{\authname}

\date{}

\input{../Comments.text}
\input{../Common.text}

\begin{document}

\maketitle

This document summarizes the contributions of each team member up to the POC
Demo.  The time period of interest is the time between the beginning of the term
and the POC demo. Project repository:
\href{https://github.com/RishabShine/IDE3.0}{https://github.com/RishabShine/IDE3.0}.

\section{Demo Plans}

\noindent\textbf{Feasibility risk.}
The primary feasibility risk is whether the team can construct and maintain a
model that accurately represents relationships between a software system's
implementation and engineering intent (requirements, constraints, and design
decisions), including justifications for those links as the code evolves. This
risk is stated in the Development Plan's Proof of Concept section
(\texttt{docs/DevelopmentPlan/DevelopmentPlan.tex}).

\noindent\textbf{What we will demonstrate.}
The PoC demo shows the foundation of that intent model in the repository's
agentic workflow under \texttt{src/agentic\_flow/}, rather than a full IDE
product surface:

\begin{itemize}
  \item \textbf{Intent capture (SSON).} The
    \texttt{capture\_intent} graph
    (\texttt{src/agentic\_flow/capture\_intent/}) elicits and evaluates a
    structured statement of need, looping until the evaluation is met.
  \item \textbf{Stakeholder identification.} The user-modelling subgraph
    (\texttt{src/agentic\_flow/user\_modelling/}) drafts and refines
    stakeholders from captured intent, with human-in-the-loop interrupts, as
    documented in
    \texttt{src/agentic\_flow/user\_modelling/README.md}.
  \item \textbf{Runnable checks.} Graph behaviour is covered by unit tests in
    \texttt{test/agentic\_flow/} (run via
    \texttt{PYTHONPATH=src python -m unittest discover -s test/agentic\_flow -v}).
\end{itemize}

Together these pieces demonstrate that the team can capture engineering intent
in a structured, revisable form and attach early project context (stakeholders)
to it---the prerequisite for later trace links from intent to implementation.
UI mock-ups of the same workflow appear in the Wireframe CD
(\texttt{docs/CDs/Wireframe/}).
\section{Team Meeting Attendance}

\begin{table}[H]
\centering
\begin{tabular}{ll}
\toprule
\textbf{Student} & \textbf{Meetings}\\
\midrule
Total & 4\\
Qusay Qadir & 4\\
Rishab Shine & 4\\
Tyler Yue & 4\\
Vedant Patel & 4\\
Saumy Patel & 4\\
\bottomrule
\end{tabular}
\end{table}

The team held a meeting on the Thursday immediately after each informal TA
document discussion. These meetings took place over the team's Discord channel
and chat rather than as separate scheduled calls. They were not originally
recorded as GitHub issues; the matching \texttt{meeting}-labelled issues were
added afterward for this deliverable so attendance could be tracked in the
repo. The four meetings in the PoC period were:

\begin{itemize}
  \item September 24, 2026 --- after Problem Statement informal (Sept~23)
  \item October 8, 2026 --- after SRS informal (Oct~7)
  \item October 29, 2026 --- after VnV Plan / PoC informal (Oct~28)
  \item November 12, 2026 --- after Design Doc Rev~-1 informal (Nov~11)
\end{itemize}

\section{Supervisor/Stakeholder Meeting Attendance}

\noindent \textbf{Supervisor's Name: } Dr.\ Istvan David

\begin{table}[H]
\centering
\begin{tabular}{ll}
\toprule
\textbf{Student} & \textbf{Meetings}\\
\midrule
Total & 4\\
Qusay Qadir & 4\\
Rishab Shine & 4\\
Tyler Yue & 4\\
Vedant Patel & 4\\
Saumy Patel & 4\\
\bottomrule
\end{tabular}
\end{table}

All five team members attended every supervisor meeting held so far in the PoC
period. Attendance was coordinated through the team's Discord channel. Matching
GitHub issues were not created after each supervisor meeting; they were added
afterward for this deliverable so the counts could be recorded in the
repository.

\section{Lecture Attendance}

\begin{table}[H]
\centering
\begin{tabular}{ll}
\toprule
\textbf{Student} & \textbf{Lectures}\\
\midrule
Total & 7\\
Qusay Qadir & 7\\
Rishab Shine & 7\\
Tyler Yue & 4\\
Vedant Patel & 1\\
Saumy Patel & 5\\
\bottomrule
\end{tabular}
\end{table}

Lecture attendance was tracked informally through the team's Discord channel
rather than through per-lecture GitHub issues at the time. The matching lecture
issues were added afterward for this deliverable so attendance could be
recorded in the repository; they were not created immediately after each
lecture.

\section{TA Document Discussion Attendance}

\noindent \textbf{TA's Name: } Lucas Dutton

\begin{table}[H]
\centering
\begin{tabular}{ll}
\toprule
\textbf{Student} & \textbf{Meetings}\\
\midrule
Total & 4\\
Qusay Qadir & 4\\
Rishab Shine & 4\\
Tyler Yue & 4\\
Vedant Patel & 0\\
Saumy Patel & 4\\
\bottomrule
\end{tabular}
\end{table}

The four informal TA document discussions in the PoC period were Problem
Statement (Sept~23), SRS (Oct~7), VnV Plan / PoC (Oct~28), and Design Doc
Rev~-1 (Nov~11). Everyone attended all four except Vedant, who attended none.
Coordination and follow-up happened primarily on Discord. Matching GitHub
issues were not created after each TA discussion; they were added afterward for
this deliverable so attendance could be recorded in the repository.

\section{Commits, Lines Added and Lines Deleted, Pull Requests}

\begin{table}[H]
\centering
\begin{tabular}{lllllll}
\toprule
\textbf{Student} & \textbf{Commits} & \textbf{Percent} & \textbf{Lines Added} &
\textbf{Lines Deleted} & \textbf{Pull Requests}\\
\midrule
Total & 175 & 100\% & 25,570 & 3,010 & 63\\
Qusay Qadir & 63 & 36.00\% & 4,183 & 1,285 & 18\\
Rishab Shine & 39 & 22.29\% & 14,582 & 249 & 14\\
Tyler Yue & 29 & 16.57\% & 2,992 & 1,023 & 13\\
Vedant Patel & 27 & 15.43\% & 2,182 & 388 & 11\\
Saumy Patel & 9 & 5.14\% & 1,631 & 65 & 7\\
\bottomrule
\end{tabular}
\end{table}

These figures were collected on November 19, 2026 from the
\href{https://github.com/RishabShine/IDE3.0}{project repository} at
\texttt{main} head \texttt{d6b654c}. The table counts all commits reachable from
\texttt{main}, including merge commits. The team total also includes 8
automated \texttt{github-actions[bot]} commits (not shown as a row), so member
percentages sum to less than 100\%. Author names were matched as follows:
Qusay Qadir; Rishab Shine (personal and GitHub no-reply emails); Tyler Yue
(personal and GitHub no-reply emails); Vedant Patel; and Saumy Patel
(\texttt{SaumyP04}). Lines added and deleted are summed from text-file changes
in non-merge commits; binary-file changes are excluded. These totals include
the initial capstone-template import and therefore do not represent only newly
written project content. Pull requests include both open and closed requests
authored by each contributor (\texttt{qusayqadir}, \texttt{RishabShine},
\texttt{Tyler-Yue}, \texttt{vp2306}, \texttt{SaumyP04}).

\section{Issue Tracker}

\begin{table}[H]
\centering
\begin{tabular}{lll}
\toprule
\textbf{Student} & \textbf{Authored (O+C)} & \textbf{Assigned (C only)}\\
\midrule
Qusay Qadir & 11 & 9 \\
Rishab Shine & 6 & 0 \\
Tyler Yue & 7 & 0 \\
Vedant Patel & 2 & 0 \\
Saumy Patel & 2 & 0 \\
\bottomrule
\end{tabular}
\end{table}

As of November 19, 2026, pull requests are excluded from these issue counts and
are counted in the commits section instead. Authored counts include open and
closed issues; assigned counts include only closed issues.

\begin{itemize}
  \item Qusay (\texttt{qusayqadir}): authored 11 issues; assigned on 9 closed
    issues.
  \item Rishab (\texttt{RishabShine}): authored
    \href{https://github.com/RishabShine/IDE3.0/issues/95}{\#95},
    \href{https://github.com/RishabShine/IDE3.0/issues/96}{\#96},
    \href{https://github.com/RishabShine/IDE3.0/issues/108}{\#108},
    \href{https://github.com/RishabShine/IDE3.0/issues/109}{\#109},
    \href{https://github.com/RishabShine/IDE3.0/issues/110}{\#110}, and
    \href{https://github.com/RishabShine/IDE3.0/issues/111}{\#111}; no closed
    issues assigned.
  \item Tyler (\texttt{Tyler-Yue}): authored
    \href{https://github.com/RishabShine/IDE3.0/issues/80}{\#80}--\href{https://github.com/RishabShine/IDE3.0/issues/82}{\#82}
    and
    \href{https://github.com/RishabShine/IDE3.0/issues/112}{\#112}--\href{https://github.com/RishabShine/IDE3.0/issues/115}{\#115};
    no closed issues assigned.
  \item Vedant (\texttt{vp2306}): authored
    \href{https://github.com/RishabShine/IDE3.0/issues/97}{\#97} and
    \href{https://github.com/RishabShine/IDE3.0/issues/98}{\#98}; no closed
    issues assigned.
  \item Saumy (\texttt{SaumyP04}): authored
    \href{https://github.com/RishabShine/IDE3.0/issues/99}{\#99} and
    \href{https://github.com/RishabShine/IDE3.0/issues/117}{\#117}; no closed
    issues assigned.
\end{itemize}

\section{CICD via GitHub Actions}

CI/CD is implemented with GitHub Actions. The active workflow builds LaTeX
documents when \texttt{docs/}, \texttt{refs/}, or related build inputs change,
commits generated PDFs back to \texttt{main} when appropriate, and deploys
project pages. That keeps documentation artifacts reproducible and reviewable
on every documentation pull request.

\begin{itemize}
  \item Workflow file:
    \href{https://github.com/RishabShine/IDE3.0/blob/main/.github/workflows/latex-pages.yml}{\texttt{.github/workflows/latex-pages.yml}}
  \item Supporting scripts and dependency manifest:
    \texttt{.github/scripts/} and \texttt{.pdf-deps.json}
\end{itemize}

Per the Development Plan, the team still intends to extend Actions to automated
tests, linting, and formatting for \texttt{src/} (for example running the
existing \texttt{test/agentic\_flow} suite on pull requests). Documentation CI
is in place now; product-code CI will be added as the implementation grows past
the PoC.

\section{Custom Documents (CDs)}

As recorded in \texttt{TeamComposition.csv}, the team's planned CDs are:

\begin{itemize}
  \item \textbf{Wireframe Report} (CD-Fall) --- delivered under
    \texttt{docs/CDs/Wireframe/}, including a clickable prototype
    (\texttt{docs/CDs/Wireframe/prototype/index.html}) aligned with the intent
    capture / stakeholder / verification workflow.
  \item \textbf{Usability Report} (later) --- planned evaluation with
    SFWRENG~2AA4 students as usability participants, per the team composition
    stakeholders.
\end{itemize}

The CD plan does not need to change: the Wireframe CD remains the right Fall
artifact for communicating the human-in-the-loop IDE surface implied by the
SRS and MG/MIS, and the Usability Report remains the right follow-on once more
of the product UI is implemented. No replacement CD types are proposed at this
time.

\section{Team Charter Trigger Items}

The Development Plan appendix (Team Charter) identifies these quantified or
operational triggers for under-contribution and process failure:

\begin{itemize}
  \item Work for each deliverable should be split into a roughly equal share of
    issues per member; progress is reviewed at team meetings.
  \item Pull requests require peer review before merge (at least one reviewer;
    the charter also aims for two reviewers on merge requests).
  \item Attendance, commit history, and issue completion are the project
    management metrics used to judge fair share.
  \item A member consistently below target is first raised with them by the
    team, then with the TA/instructor if it continues.
  \item Members must flag blockers early rather than missing internal deadlines
    silently.
\end{itemize}

\noindent\textbf{Violations (aggregate).}
In the PoC period, informal TA document discussions were missed by one member
(0 of 4). Lecture attendance is uneven across the team. Team and supervisor
meeting attendance were complete. No member was escalated to the TA or
instructor for under-contribution. Meeting issues were created late (after the
fact for this report), so attendance triggers were harder to audit in-repo
during the term.

\noindent\textbf{Plan.}
\begin{itemize}
  \item Create meeting, lecture, and TA issues at the time of each event going
    forward so triggers are measurable from the repository.
  \item Keep the equal-issue-share and PR-review rules; treat missed informal
    TA sessions as a soft trigger that requires a written Discord debrief from
    the absent member within 24 hours.
  \item Revisit lecture-attendance expectations after PoC: either record
    lecture issues consistently or narrow the trigger to deliverable-critical
    sessions only, so the metric matches how the team actually works.
\end{itemize}

\section{Additional Productivity Metrics}

Beyond the attendance, commit, PR, and issue tables above, the team tracked:

\begin{itemize}
  \item \textbf{Pull requests authored:} 63 total on the repository (see
    Section~6), used as a proxy for branch-based delivery.
  \item \textbf{Automated graph tests:} the agentic-flow unit suite under
    \texttt{test/agentic\_flow/} (stakeholder graph coverage with controlled
    model responses; no live Ollama required).
  \item \textbf{Documentation CI runs:} LaTeX/PDF workflow executions on
    documentation pull requests via
    \texttt{.github/workflows/latex-pages.yml}.
\end{itemize}

The team is not currently counting late arrivals, correction cycles, or missed
internal deadlines as separate numeric metrics; those will be added only if
charter triggers continue to be hard to interpret from GitHub data alone.

\end{document}
